\documentclass[11pt]{article}
\usepackage{amsmath, amssymb}
\usepackage{geometry}
\usepackage{listings}
\usepackage{xcolor}
\usepackage{hyperref}
\geometry{margin=2.3cm}

\lstset{
    basicstyle=\ttfamily\small,
    columns=fullflexible,
    keepspaces=true,
    breaklines=true,
    frame=single,
    framerule=0.3pt,
    rulecolor=\color{gray!50},
    backgroundcolor=\color{gray!5},
    aboveskip=4pt,
    belowskip=4pt,
}

\newcommand{\file}[1]{\texttt{#1}}
\newcommand{\code}[1]{\texttt{#1}}
\newcommand{\dd}{\mathrm{d}}
\newcommand{\Tr}{\mathrm{Tr}}
\newcommand{\calN}{\mathcal{N}}

\title{Summary of changes: pCN sampler for \texttt{spinDMFT}}
\date{}

\begin{document}
\maketitle

\section{Overview}

The original \file{spinDMFT/Algorithm/} samples the mean-field trajectory $V$ from
the Gaussian prior $p(V)$ and reweights observables by the per-trajectory
partition function $Z(V) = \Tr[U(\beta,0;V)]$.  At low temperatures $Z(V)$ becomes
sharply peaked in the prior support, leaving the importance-sampling estimator
weight-degenerate and forcing extremely large sample counts.

The new directory \file{spinDMFT/Algorithm\_MH/} is a self-contained copy of the
algorithm that replaces importance sampling with a preconditioned Crank--Nicolson
(pCN) Markov chain~\cite{Cotter2013} targeting
\begin{equation}
    \pi(\dd V) \propto p(V)\,Z(V)\,\dd V,
\end{equation}
the measure under which the sample mean of $\Tr[U(\beta{-}\tau)\,S^\alpha\,U(\tau)\,S^\beta]/Z(V)$
is already the thermal correlation $g^{\alpha\beta}(\tau)$.  No reweighting is
needed; for Gaussian-prior, positive-likelihood targets pCN has dimension-independent
autocorrelation, so the effective sample size is on par with that of independent
draws from the target.

\section{File-by-file changes}

\subsection{\file{Parameter\_Space/}}

Three new CLI options are exposed:

\begin{itemize}
    \item \code{--mhStepSize} (\code{RealType}, default $0.3$).  The pCN step
          parameter $\beta_{\mathrm{pCN}}\in(0,1]$.  Larger $\beta_{\mathrm{pCN}}$
          drives the proposal closer to an independent draw from the prior; smaller
          $\beta_{\mathrm{pCN}}$ takes more local moves.
    \item \code{--mhBurnIn} (\code{size\_t}, default $100$).  Number of cold-start
          warm-up pCN steps per MPI core, used in the first self-consistent iteration
          and discarded before observables are accumulated.
    \item \code{--mhWarmBurnIn} (\code{RealType}, default $0.25$).  Fraction of
          \code{mhBurnIn} used for the warm-start refresh in every self-consistent
          iteration after the first.  At iteration $k>1$ the chain is seeded with the
          previous iteration's final $V$ rotated into the new diagonal basis, so a
          fraction of the cold burn-in length is typically sufficient.  Set to $0$
          to skip the refresh entirely (maximum speedup, but watch for residual bias
          when early-iteration acceptance is low); set to $1$ to recover the old
          per-iteration cold burn-in cost.
\end{itemize}

The corresponding members \code{ParameterSpace::mh\_step\_size},
\code{ParameterSpace::mh\_burn\_in}, and
\code{ParameterSpace::mh\_warm\_burn\_in\_frac} are added to the header.

\subsection{\file{Run\_Time\_Data/}}

A small acceptance-rate diagnostic is added.  Per MPI core, the chain counts
proposed and accepted moves into \code{mh\_proposed\_count} and
\code{mh\_accepted\_count}; at the end of each self-consistent iteration the new
method \code{record\_mh\_acceptance()} performs an
\code{MPI\_Allreduce}\,$(\mathrm{MPI\_SUM})$ and appends the global ratio to the
member \code{acceptance\_rates}, which is also printed at iteration end and
persisted to the HDF5 output (\S\ref{sec:diag}).

Several IS-only accumulators are removed since the MH chain produces unweighted
samples from $\pi$: \code{sample\_cov\_Re}, \code{sample\_cov\_Im},
\code{spin\_expval\_cov}, and \code{Z\_sqsum} are gone.  The standard-error
computation in \code{compute\_sample\_stds} is correspondingly rewritten as the
textbook formula
\begin{equation}
    \mathrm{stderr}(\langle o\rangle)
    = \sqrt{ \bigl( \langle o^2\rangle - \langle o\rangle^2 \bigr) / N },
\end{equation}
where $\langle o\rangle$ is the (already normalized) Monte-Carlo mean,
$\langle o^2\rangle = \sum_i o_i^2 / N$ comes from the running sum-of-squares
\code{sample\_sqsum\_*}, and $N$ is the total number of MC samples across all
cores.  Both correlations and spin expectation values share this single helper.

\subsection{\file{Functions/Functions.h} and \file{Functions.cpp}}

Two new pieces live in \code{namespace spinDMFT::Functions}:

\begin{itemize}
    \item \code{diag\_freq\_to\_time(V\_diag, ortho)} -- given the chain state
          expressed in the diagonal basis of the per-block prior covariance
          (eigenvalues $\lambda_{n,\alpha}$), returns the corresponding
          imaginary-time trajectory $V(\tau_n)$.  Mirrors the back-transform
          conventions of
          \code{fmvg::FrequencyNoiseVectors::fourier\_back\_transform}.
    \item Class \code{PCNChain} -- the entire pCN chain, described in
          \S\ref{sec:pcnchain}.
\end{itemize}

\textbf{The single most important change to the existing helpers} is to
\code{compute\_spin\_correlations}.  The original IS estimator accumulates the
unweighted trace
$\Tr[U(\beta{-}\tau)\,S^\alpha\,U(\tau)\,S^\beta]$ and the partition function $Z$
separately, dividing at the end.  Under MH sampling from $\pi$, the per-sample
observable is already the correlation
\begin{equation}
    o^{\alpha\beta}(\tau; V) =
    \frac{\Tr[U(\beta{-}\tau;V)\,S^\alpha\,U(\tau;V)\,S^\beta]}{Z(V)},
\end{equation}
so each per-sample contribution is divided by $Z(V)$ inside the accumulator.
The function now reads (schematically)
\begin{lstlisting}
const RealType inv_Z = ( Z != 0 ) ? 1. / Z : 0.;
const RealType o     = std::real( single_sample_result ) * inv_Z;
spin_correlations    += o;     // sum of o_i, divided by N at the end -> <o>
sample_sqsum         += o * o; // sum of o_i^2, used in stderr formula
\end{lstlisting}
Because the chain produces unweighted samples from $\pi$, there is no IS
covariance term: \code{MPI\_share\_results} and \code{compute\_sample\_stds} drop
the \code{sample\_cov\_*} and \code{Z\_sqsum} machinery entirely (see the
\file{Run\_Time\_Data/} section).  The global ``partition function'' passed to
\code{normalize} is replaced by the total sample count
$N = \texttt{my\_pspace.num\_Samples}$.  Several dead parameters were also
dropped along the way: \code{compute\_spin\_correlations} no longer takes
\code{pspace}, and \code{normalize} no longer takes \code{rtdata}.  The MPI
helper \code{MPI\_share\_results} grew a small private utility
\code{mpi\_sum\_fieldvector} that calls \code{MPI\_Allreduce} directly on the
contiguous storage of \code{FieldVector}, removing a few rounds of buffer
copying.

\subsection{\code{PCNChain} class (lives in \file{Functions/})}\label{sec:pcnchain}

All of the chain machinery is encapsulated in a single class whose only public
surface is a constructor, a \code{step()} method, and a handful of accessors.

\begin{itemize}
    \item \textbf{State held by the chain.}  Diagonal-basis position
          $V_{\mathrm{diag}}$, current $Z(V)$, the cached propagators
          $U(\tau)$ and $U(\tau)^{-1}$, and the cached spin trajectories
          $S^\alpha(\tau)$.  Also held: a per-mode \code{std::normal\_distribution}
          for proposal noise, and a $\mathrm{Uniform}(0,1)$ for accept/reject.

    \item \textbf{Cold-start constructor.}  Builds the per-mode prior
          distributions from the (possibly truncated) eigenvalues, draws an initial
          $V_{\mathrm{diag}}$ from the prior, and retries up to $32$ times if
          $Z(V)\le 0$.

    \item \textbf{Warm-start constructor.}  Same construction, but seeds
          $V_{\mathrm{diag}}$ from a caller-supplied vector instead of the prior
          draw.  Components in modes that the new prior has frozen
          ($\lambda\le 0$) are forced to zero.  Used to carry chain state across
          self-consistent iterations (\S\ref{sec:warm-start}).

    \item \textbf{\code{step()}.}  Proposes
          $V'_{\mathrm{diag}} = \sqrt{1-\beta_{\mathrm{pCN}}^2}\,V_{\mathrm{diag}}
          + \beta_{\mathrm{pCN}}\,\xi$ with $\xi\sim p$, computes $Z(V')$,
          accepts with probability $\min(1, Z'/Z)$ ($Z\le 0$ treated as
          forbidden), and on acceptance refreshes the cached propagators and
          $S^\alpha(\tau)$.

    \item \textbf{Accessors.}  \code{Z()}, \code{S\_x()}, \code{S\_y()},
          \code{S\_z()} for the observable accumulator; \code{V\_full()} returns
          the chain state in the basis-independent Matsubara representation
          $V_{\mathrm{full}}[b] = O[b]\,V_{\mathrm{diag}}[b]$ for cross-iteration
          state transfer.
\end{itemize}

\subsection{\file{main.cpp}}\label{sec:warm-start}

The Monte-Carlo loop becomes a thin shell around \code{PCNChain}.  Across the
self-consistent iterations the loop carries a per-rank \code{std::vector<fmvg::Vector>
my\_V\_full\_persistent} that is empty before the first iteration and holds the
previous iteration's final $V_{\mathrm{full}}$ thereafter.

\begin{enumerate}
    \item \textbf{Chain construction.}  If \code{my\_V\_full\_persistent} is
          empty (first iteration), build the chain with the cold-start
          constructor.  Otherwise rotate the saved $V_{\mathrm{full}}$ into the
          new diagonal basis,
          \begin{equation}
              V_{\mathrm{diag, new}}[b] = O_{\mathrm{new}}[b]^{\top}\,
              V_{\mathrm{full}}[b],
          \end{equation}
          and pass it to the warm-start constructor.

    \item \textbf{Burn-in.}  In the first iteration, run \code{mhBurnIn} pCN
          steps.  In every subsequent iteration, run
          $\bigl\lfloor \text{mhWarmBurnIn} \cdot \text{mhBurnIn}\bigr\rfloor$
          refresh steps.  Acceptance counters update in both phases for
          diagnostics.

    \item \textbf{Production.}  \code{num\_SamplesPerCore} pCN steps; each step
          calls \code{chain.step()} and feeds the (cached or newly rebuilt)
          $(\text{chain}.Z(), \text{chain}.S_x(), \dots)$ to
          \code{compute\_spin\_correlations}.  On a rejected step the same
          trajectory is re-accumulated, which correctly weights the current
          state by its holding time.

    \item \textbf{State capture.}  After the production sweep,
          \code{my\_V\_full\_persistent = chain.V\_full()} stores the chain state
          in the basis-independent representation, ready for the next
          iteration's warm rotation.

    \item \textbf{Normalization plumbing.}  The IS variable
          \code{my\_partition\_function} is gone; in its place
          $N = \texttt{my\_pspace.num\_Samples}$ is passed directly to
          \code{normalize} and \code{compute\_sample\_stds}.  After
          \code{record\_mh\_acceptance()}, the iteration finishes through the
          existing convergence check and storage path unchanged.
\end{enumerate}

\paragraph{Empirical note on warm-start.}  When early-iteration acceptance is low
($\lesssim 20\%$, as happens at low temperature before the self-consistent
covariance settles), a small \code{mhWarmBurnIn} fraction can carry residual
bias from the initial chain state into the final estimator.  At
$\beta J = 12$, \code{mhBurnIn} $= 2000$, \code{mhWarmBurnIn} $= 0$ gives a
$\sim 5\sigma$ bias relative to the cold-only result; \code{mhWarmBurnIn} $= 1$
recovers it.  The default of $0.25$ is a safe choice for high-temperature
sweeps; raise it towards $1$ in deep-cold regimes.

\subsection{\file{CMakeLists.txt}}

The output binary is renamed from \code{executable\_DOUBLE.out} to
\code{executable\_MH\_DOUBLE.out} so the pCN binary and the original IS binary
can coexist in \file{spinDMFT/}.

\section{Mathematical contract of the new estimator}

The original IS code computes
\begin{equation}
    \widehat{g}_{\mathrm{IS}}^{\,\alpha\beta}(\tau)
    = \frac{\sum_{i=1}^N \Tr[U(\beta{-}\tau;V_i)\,S^\alpha\,U(\tau;V_i)\,S^\beta]}
           {\sum_{i=1}^N Z(V_i)},
    \qquad V_i \sim p(V).
\end{equation}
The new code computes
\begin{equation}
    \widehat{g}_{\mathrm{pCN}}^{\,\alpha\beta}(\tau)
    = \frac{1}{N}\sum_{i=1}^N
      \frac{\Tr[U(\beta{-}\tau;V_i)\,S^\alpha\,U(\tau;V_i)\,S^\beta]}{Z(V_i)},
    \qquad V_i \sim \pi.
\end{equation}
Both target the same thermal expectation $g^{\alpha\beta}(\tau)$.  The pCN form
avoids the IS denominator, which is the source of the weight-degeneracy problem
at low temperature: when $Z(V)$ is sharply peaked the IS estimator concentrates
all its statistical weight on a vanishing fraction of the prior samples, whereas
the pCN chain draws preferentially from the peaked region and weights every
sample equally.

\section{Diagnostic interface}\label{sec:diag}

Every self-consistent iteration now prints

\begin{center}
\code{MH acceptance rate            : <number>}
\end{center}

after the standard convergence-error report.  The full per-iteration time series
is held in \code{RunTimeData::acceptance\_rates} during the run and persisted to
the HDF5 output as the attribute
\code{runtimedata/mh\_acceptance\_rates} (a list with one entry per completed
self-consistent iteration), alongside the existing
\code{absolute\_iteration\_errors} and \code{negative\_eigenvalue\_ratios}.  A
practical tuning target is $30$--$60\%$ acceptance; below $\sim 20\%$ shrink
\code{--mhStepSize}, above $\sim 80\%$ enlarge it.

\section{Tested regime}

Convergence has been verified against Grempel's reference (\file{Data/grempel\_plot.csv})
at $\beta J \in \{1, 2, 3, 5, 7, 12, 30\}$ with \code{numTimeSteps=100} and
\code{cstype=A}.  Relative agreement at $\tau=\beta/2$ is everywhere
$\le 1.2\%$.  At $\beta J = 30$ the chain converges in $\sim 5$ self-consistent
iterations with \code{--mhStepSize=0.3} (acceptance $\sim 50\%$); longer chains
($\sim 8 \times 10^5$ total samples, $8\times 10^4$ burn-in) bring the
$\tau \leftrightarrow \beta{-}\tau$ asymmetry below $4\times 10^{-4}$.

\begin{thebibliography}{1}
\bibitem{Cotter2013}
S.~L.~Cotter, G.~O.~Roberts, A.~M.~Stuart, D.~White,
\textit{MCMC methods for functions: modifying old algorithms to make them faster},
Statistical Science \textbf{28}, 424 (2013).
\end{thebibliography}

\end{document}
